% Generated draft: requires mathematical and source review.
\documentclass[11pt]{article}
\usepackage[margin=25mm]{geometry}
\usepackage{fontspec}
\setmainfont{Latin Modern Roman}
\usepackage{amsmath,amssymb,mathtools,unicode-math}
\setmathfont{Latin Modern Math}
\usepackage{enumitem,needspace,xurl,hyperref}
\setlength{\parindent}{0pt}
\setlength{\parskip}{5pt}
\setlength{\emergencystretch}{3em}
\newcommand{\N}{\mathbb{N}}
\newcommand{\Z}{\mathbb{Z}}
\newcommand{\Q}{\mathbb{Q}}
\newcommand{\R}{\mathbb{R}}
\newcommand{\C}{\mathbb{C}}
\newcommand{\F}{\mathbb{F}}
\newcommand{\A}{\mathbb{A}}
\newcommand{\PP}{\mathbb{P}}
\newcommand{\E}{\mathbb{E}}
\newcommand{\eps}{\varepsilon}
\newcommand{\dd}{\,\mathrm{d}}
\begin{document}
\section*{UnsolvedMath ID: 800823}
\textbf{Source problem code:} catalogue entry
\subsection{Definitions and mathematical statement}
Let \(M\) denote the intended two-dimensional fluid domain or surface. Its geometry, topology, metric, and boundary conditions are not specified in the supplied record. In an appropriate setting, write the incompressible Euler equations schematically as
\[
\partial_t u+\nabla_u u=-\nabla p,\qquad \operatorname{div}u=0,
\]
where \(u\) is the fluid velocity and \(p\) is the pressure. With scalar vorticity \(\omega=\operatorname{curl}u\), the corresponding vorticity equation is
\[
\partial_t\omega+u\cdot\nabla\omega=0.
\]
The relation recovering \(u\) from \(\omega\), including any circulation data, depends on the unspecified setting.

A point-vortex configuration is represented formally by a vorticity measure
\[
\omega_{\mathrm{pv}}(t)=\sum_{j=1}^{N}\Gamma_j\delta_{z_j(t)},
\]
where \(N\) is a positive integer, \(\Gamma_j\in\R\) are vortex strengths, \(z_j(t)\in M\) are vortex positions, and \(\delta_z\) denotes the Dirac measure at \(z\). The associated finite-dimensional dynamics, its collision restrictions, and the meaning of integrability depend on the chosen fluid setting. The supplied record does not identify the intended integrable cases or specify whether an additional background vorticity is present.

\textbf{Problem (Khesin, Misiolek, and Modin, 2022, Problem 22).} Justify the phenomenon reported in the references labelled \([\mathrm{ModViv19}]\) and \([\mathrm{ModViv21}]\): integrable cases of point-vortex dynamics appear to act as attractors for two-dimensional Euler flows with generic smooth initial vorticity when small numerical viscosity is present. In addition, design a model of dissipation in the two-dimensional Euler equation that produces such integrable dynamics at large times.

\textbf{Unresolved formulation.} This is a qualitative research problem, rather than a fully quantified conjecture in the supplied record. A precise formulation requires specifying \(M\), boundary and circulation conditions, a space of smooth initial vorticities and a meaning of ``generic,'' the relevant integrable point-vortex systems, and the observables or topology in which attraction is asserted. It also requires specifying how smooth vorticity is related to the proposed vortex dynamics, the numerical or continuous dissipation model, and the order or relationship of the large-time and small-viscosity limits. The record does not assert that the inviscid Euler equation itself has the proposed attracting behavior, nor does it supply a rigorous attraction theorem.
\subsection{Short English statement}
Explain the reported numerical tendency of two-dimensional Euler flows with generic smooth initial vorticity and small numerical viscosity to approach integrable point-vortex behavior. Construct a dissipative Euler model exhibiting this behavior at large times. The domain, integrable cases, genericity, attraction criterion, and limiting procedure remain unspecified.
\subsection{Sources}
\begin{itemize}
\item Supplied local UnsolvedMath registry record: permanent ID \(800823\), exact source problem code \texttt{catalogue entry}, titled ``Geometric Hydrodynamics in Open Problems --- Problem 22.'' The registry attributes the problem to Khesin, Misiolek, and Modin (2022) and labels it \texttt{open}; this is supplied status metadata, not an independently verified mathematical assessment.
\item Khesin, Misiolek, and Modin, \emph{Geometric Hydrodynamics in Open Problems} (2022), Problem 22. Supplied source URL: \url{https://arxiv.org/abs/2205.01143}. The background records access on 2026-10-16 and extraction by \texttt{arxiv\_latex}; these are supplied provenance claims, not access or verification performed for this draft.
\item The statement attributes the observed phenomenon to references labelled \([\mathrm{ModViv19}]\) and \([\mathrm{ModViv21}]\). Their bibliographic entries and links are absent from the supplied record, so their identities and precise results remain unresolved.
\end{itemize}
\end{document}
